% s-339-p339-m002; original PDF 339, 340
\Aphorism{s-339-p339-m002}{389}
Statt eine objective Natur anzusetzen und durch Reaction gegen diese das menschliche Affect- und Geistesleben entstehen zu lassen (Anpassung an äussere Bedingungen, nach Spencer), kann man vorziehen, umgekehrt dem Triebleben des Menschen die erste Stelle einzuräumen und daraus durch spontane Acte (Interpretation, Aneignung, Wille zur Macht) eine objective Natur, ein Weltbild, herauswachsend zu denken, als einen langsam sich ansammelnden Bestand von Interpretationen, über deren Richtung sich verschiedene Individuen geeinigt haben. Um sich dem Nachbar mitzutheilen, — eine Nothwendigkeit, die in jeder biologischen Reihe sehr früh eintritt — verzichtet man auf seine volle persönliche Willkür im Auffassen, Anschauen, Deuten, Namengeben; was auf diese Weise an Gleichartigkeit erreicht und erblich festgehalten wird, scheint dann später den Charakter des Willkürlichen ganz eingebüsst zu haben und gilt als das Gegentheil des Beliebens, als Object, Wissen, Erkenntniss. — Ist also der Mensch eine Reaction gegen Natur oder Natur eine Action des Menschen? Vielleicht ist es am räthlichsten, beide Formeln neben einander stehen zu lassen; wenn man zu diesem Scheingegensatz von These und Antithese die Synthese finden sollte, so wäre damit auch der Stein der Weisen, nämlich die ideale Erkenntnisstheorie, gefunden.
\par

